% Imporditud: tools/import_eksperimendid.py
% Allikas: Eesti füüsikaolümpiaadi eksperimendi ülesannete
%   kogu (Rael Kalda, 2023), ülesanne Ü65, lahendus L65.
\setAuthor{Erkki Tempel}
\setRound{lõppvoor}
\setYear{2015}
\setNumber{P E1}
\setDifficulty{2}
\setTopic{vedelike mehaanika}
\setVahendid{U-toru, millimeeterpaber, vesi ($\rho$ = 1,0 g/cm$^3$), piiritus, toiduõli,}
\setPunktid{10}
\setAllikas{Eksperimendi kogu Ü65, lahendus lk 61}
\setLahendusseis{toores}

\prob{U-toru}
Leidke piirituse tihedus.

statiiv käpaga.

\hint

\solu
Täidame U-toru poolenisti veega ning lisame ühte U-toru

harusse õhukese kihi toiduõli. Toiduõli lisamine on vaja-

lik, et vesi ja piiritus ei saaks omavahel seguneda Nüüd lisa- vesi piiritus me ettevaatlikult õli peale piiritust. Piirituse samba kõrgus $h_2$

peaks olema vähemalt 5 cm. Kuna toiduõli kiht on võrreldes $h_1$

piirituse kihiga väga väike, siis võib selle jätta arvestamata

(lisada piirituse nivoo juurde).

õli

pp = pv $\Rightarrow$ $\rho$pgh2 = $\rho$vgh1 $\Rightarrow$ $\rho$p = $\rho$v $h_1$ $\approx$ 0,8 g/cm$^3$ vesi $h_2$
\probend
